% ICASSP draft -- numerical claims marked TODO must be replaced/confirmed by
% additional experiments. Results not marked TODO are reproduced from our own
% internal experiment log (single run per setting unless noted); error bars
% and multi-seed averages are still required before submission.
\documentclass{article}
\usepackage{spconf,amsmath,amssymb,graphicx,booktabs,algorithm,algorithmic,hyperref}

\newcommand{\x}{\mathbf{x}}
\newcommand{\g}{\mathbf{g}}
\newcommand{\C}{\mathcal{C}}
\newcommand{\Sset}{\mathcal{S}}
\newcommand{\R}{\mathbb{R}}
\newcommand{\todo}[1]{\textbf{[TODO: #1]}}

\title{ONE LUCKY DRAW IS ENOUGH: STOCHASTIC GRADIENT LEAKAGE UNDER\\
QUANTIZATION AND SPARSIFICATION, AND A CONSISTENCY-RANKED RESTART ATTACK}
\name{Anonymous Author(s)}
\address{Anonymous Affiliation}

\begin{document}
\maketitle

\begin{abstract}
This paper proposes CoRA (Consistency-Ranked restart Attack), a restart-based
gradient inversion attack that treats communication compression in federated
learning (FL) as a probabilistic rather than absolute privacy barrier.
Quantization and sparsification are widely reported to defeat attacks such as
DLG and the Cocktail Party Attack (CPA), but this conclusion rests on
single-trial evaluations. Unlike single-shot attacks that decode the compressed
update deterministically, CoRA resamples the information destroyed by
compression from a calibrated per-tensor model and selects among candidate
reconstructions by a self-supervised compression-consistency score rather than
averaging them. We investigate how reconstruction quality changes with quantization
bit-width and sparsification retention ratio. Our goal is to determine
how far CoRA remains effective as compression becomes more aggressive,
and to identify its reconstruction limits under a specified attack budget.

\end{abstract}

\begin{keywords}
gradient inversion, federated learning, privacy, quantization, sparsification
\end{keywords}

\section{Introduction}
Federated learning (FL) keeps raw training examples on-device and
communicates only gradients or model updates. Gradient inversion attacks show
that this is not sufficient for privacy: optimization-based attacks such as
Deep Leakage from Gradients (DLG) \cite{zhu2019deep} and Inverting Gradients
\cite{geiping2020inverting}, and analytical attacks such as the Cocktail
Party Attack (CPA), which treats fully-connected gradients as a blind source
separation problem solved with independent component analysis (ICA)
\cite{kariyappa2023cocktail}, can reconstruct client inputs from communicated
gradients with no auxiliary data and no change to training.

Because these attacks were originally demonstrated on full-precision, dense
gradients, a natural question is whether the lossy communication-compression
techniques that FL systems already use for bandwidth reasons---fixed-bit
quantization of model updates and top-$k$ gradient sparsification---also
double as a privacy defense. Prior evaluations, including our own
reproduction below, generally answer yes: a single deterministic attack
attempt against a 4-bit quantized or aggressively sparsified update
reconstructs images that are visually and quantitatively far worse than the
full-precision baseline, and this is often reported as evidence that
compression ``defeats'' the attack.

We argue that this single-trial framing understates the actual risk. The
information destroyed by quantization or sparsification is destroyed
\emph{stochastically} with respect to the adversary: the client's true
communicated bits determine only an ambiguity set of gradients (and, after
inversion, of candidate images) consistent with the observation, and a fixed
deterministic decoding rule (e.g., naive dequantization, or zero-filling
pruned coordinates) picks one arbitrary point of that set. Nothing prevents
an adversary from resampling within the ambiguity set and re-running the
attack. We measured exactly this: reconstructing a sparsified update with
independently re-imputed missing coordinates and re-running CPA/ICA succeeds
in only about 4\% of independent trials in our setup, yet 4\% is not zero,
and trials are cheap and embarrassingly parallel for an honest-but-curious
server with modest compute.

This paper makes three contributions.
\begin{itemize}
  \item We empirically characterize reconstruction quality under
  quantization and sparsification, including an 8-to-4-bit
  randomized-decoding sweep and a matched sparsification-retention
  sweep. We examine how reconstruction quality changes with
  compression severity, distinguishing measured degradation from
  a calibrated success/failure boundary (Sections~3 and~5).

  \item We show that naive random resampling followed by averaging
  does not necessarily improve image reconstruction. In the tested
  4-bit setting, ensemble averaging improves recovered-source cosine
  similarity relative to a single randomized trial, but does not
  improve the reported image metrics and remains worse than naive
  deterministic decoding. This highlights the distinction between
  source-recovery quality and final image-reconstruction quality
  (Section~3.3).

  \item We propose CoRA, a restart-based framework that generates
  candidate updates by random sampling within compression-implied
  bounds and replaces averaging with attacker-observable candidate
  selection. We study feature-consistency screening with a total
  variation term for quantized updates and recompressed-update
  consistency ranking for sparsified updates, without using private
  image metrics for selection (Section~4).
\end{itemize}

\todo{Define and validate an image-recovery success criterion before
claiming specific successful and failing bit widths or retention ratios.}

\todo{Use controlled ablations to determine whether sampling-model
mismatch or candidate averaging explains the observed degradation.}

\todo{Complete the matched randomized sparsification sweep and evaluate
CoRA against deterministic decoding and averaging under equal total
compute budgets. Report measured gains or limitations without assuming
that randomized selection is superior.}

\section{Background and threat model}
\subsection{Federated update compression}
We consider FedAvg-style cross-device FL in which a client sends a local
model update $\g\in\R^{d}$ to an honest-but-curious server. Two
communication-compression schemes are in scope. \emph{Quantization} applies
fixed-bit symmetric uniform quantization to $\g$; the server observes the
dequantized update $\C_b(\g)$, with quantization bin width $\Delta$ set by
the bit-width $b$ and the per-tensor dynamic range. \emph{Sparsification} retains only the top-$k$ coordinates of $\g$
by magnitude within the specified compression scope, and the server
observes $\C_k(\g)$, an update with the pruned coordinates set to zero
and their original indices known.


\subsection{Adversary}
The server knows the model architecture, loss, and compression scheme
(bit-width or retention ratio $k$) but not the values removed by compression.
It observes one client's compressed update from one or a few local steps,
which is the standard analytical-attack setting used by DLG-style and CPA-style
evaluations \cite{zhu2019deep,kariyappa2023cocktail}, and can run any number
of independent attack attempts offline without further contacting the client.
The server does not need an auxiliary dataset or a pretrained generative
prior for the attack we study, in contrast to attacks that strengthen the
optimization prior itself \cite{li2022ggl}.

\section{Compression defeats the average attack, not every attack}
\subsection{Single-trial results reproduce the standard picture}
We first reproduce the standard single-trial picture using two setups.
\textbf{Setup A (single-client, single-image):} Tiny-ImageNet, a 2-layer
fully connected network (FC-2), Adam training (lr $10^{-3}$, batch size 128),
and a CPA-style reconstruction objective combining ICA source-separation
terms (negentropy, decorrelation) with a total-variation image prior,
optimized for 25000 iterations. \textbf{Setup B (batched, CPA/FastICA):}
ImageNet with a VGG16 client model, FedAvg local updates as the attack
input, CPA/FastICA batch size 128, 25000 CPA iterations and 25000 feature-inversion
iterations per round, following the CPA protocol of \cite{kariyappa2023cocktail}.

\begin{table}[t]
\caption{Setup A: single-trial reconstruction under compression (FC-2,
Tiny-ImageNet, 25000 attack iterations). Global refer to how
top-$k$ sparsification selects retained coordinates.}
\label{tab:fc2}
\centering
\small
\begin{tabular}{lccc}
\toprule
Setting & PSNR$\uparrow$ & SSIM$\uparrow$ & LPIPS$\downarrow$\\
\midrule
FP32 (no compression) & 26.44 & 0.794 & 0.243\\
8-bit quantization & 26.24 & 0.779 & 0.257\\
4-bit quantization & 21.82 & 0.164 & 0.702\\
Top-$k$ 80\% (global) & 25.56 & 0.714 & 0.313\\
Top-$k$ 50\% (global) & 22.92 & 0.334 & 0.602\\
Top-$k$ 10\% (global) & 21.99 & 0.197 & 0.696\\
Top-$k$ 80\% (layerwise) & 12.16 & 0.303 & 0.767\\
Top-$k$ 50\% (layerwise) & 12.16 & 0.305 & 0.769\\
Top-$k$ 10\% (layerwise) & 21.53 & 0.263 & 0.632\\
\bottomrule
\end{tabular}
\end{table}

Table~\ref{tab:fc2} shows a clear cliff: 8-bit quantization and 80\% top-$k$
retention (global or layerwise) leave reconstruction close to the FP32
baseline, while 4-bit quantization and 50\% or 10\% retention collapse SSIM
toward the level of an unstructured guess. Figure~\ref{fig:cliff} visualizes
the same cliff. A single deterministic attack attempt would conclude that
4-bit quantization and $\le$50\% sparsification are effective defenses.

\begin{figure}[t]
  \centering
  \includegraphics[width=\columnwidth]{figures/compression_cliff}
  \caption{Single-trial PSNR and SSIM of the reproduced attack (Setup A)
  across compression settings. Blue bars sit above an illustrative
  ``reconstruction succeeds'' line; gray bars sit below it. This single-trial
  view is what a naive defense evaluation would report.}
  \label{fig:cliff}
\end{figure}

\subsection{A single attempt is not the adversary's only option}
A deterministic decoder collapses the compression ambiguity set to one
point: naive dequantization maps every value in a quantization bin to its
bin center, and naive sparsification decoding maps every pruned coordinate
to zero. Neither choice is privileged from the adversary's perspective---both
are simply one sample from the set of updates consistent with the
observation. We therefore repeated the attack under Setup B with
\emph{randomized} decoding: for sparsification, each pruned coordinate is
redrawn per trial within the range implied by the top-$k$ threshold (missing
coordinates only; observed coordinates are kept fixed), and the resulting
per-trial CPA/FastICA reconstructions are treated as an ensemble of
independent attempts rather than a single answer.

\begin{table}[t]
\caption{Setup B: CPA/FastICA on VGG16/ImageNet FedAvg updates under
4-bit quantization, comparing naive and randomized decoding. ICA/source CS
is the cosine similarity of recovered vs. true mixing sources.}
\label{tab:cpa}
\centering
\scriptsize
\begin{tabular}{lcccc}
\toprule
Method & CS$\uparrow$ & PSNR$\uparrow$ & SSIM$\uparrow$ & LPIPS$\downarrow$\\
\midrule
FP32 (no compression) & 0.987 & 22.40 & 0.183 & 0.521\\
Naive int4 dequant. & 0.691 & 12.22 & 0.329 & 0.671\\
Randomized, 1 trial & 0.135 & 11.30 & 0.106 & 0.737\\
Randomized, $K{=}8$ avg. & 0.809 & 13.22 & 0.337 & 0.737\\
\bottomrule
\end{tabular}
\end{table}

Randomized decoding allows the adversary to explore multiple candidate
completions of the same compressed observation. Its effectiveness,
however, depends on the decoding and reconstruction procedure.
In a separate 8-bit quantization experiment, the randomized ensemble
improved the mean reconstruction metrics relative to naive
dequantization, whereas the 4-bit results in Table~\ref{tab:cpa}
showed no corresponding improvement. These experiments used different
protocols and therefore do not isolate the effect of bit width.
A single deterministic reconstruction does not characterize the
variation across alternative decodings, while an averaged ensemble
does not establish the benefit of retaining and evaluating individual
attempts.

\subsection{Why naive multi-trial averaging still under-delivers}
Table~\ref{tab:cpa} also shows that the obvious way to turn independent
trials into a better attack---resample and average---does not work well
here. Averaging $K{=}10$ randomized reconstructions raises the ICA/source
cosine similarity (0.135 $\to$ 0.441) but leaves PSNR/SSIM essentially
unchanged and still far below the naive int4 baseline. We attribute this to
two compounding issues, both directly observable in our logs. First, the
imputation noise was drawn i.i.d.\ per coordinate from a fixed range, but the
true information removed by compression is far from i.i.d.: each parameter
tensor has a different scale and a strongly non-uniform coefficient
distribution, so an unstructured noise model perturbs the true structure of
the update rather than recovering it. Second, pixel-space (or source-space)
averaging blends reconstructions of varying quality together; a few good
draws are diluted by many mediocre ones instead of being identified and kept.
Enforcing quantization consistency as a hard constraint inside the
reconstruction optimizer---which we also tested, both at the ICA stage and
end-to-end through feature inversion---did not fix this either: the
compressed-gradient consistency term makes the optimization landscape
coarser and more constrained without proportionally improving reconstruction
quality. Section~4 turns these three negative findings into design
requirements.

\subsection{A closed-form model of cumulative compromise}
Model each independent attack trial $t=1,\dots,K$ against a fixed client and
a fixed observed compressed update as an independent Bernoulli success
indicator $Y_t\sim\mathrm{Bernoulli}(p)$, where independence follows from
redrawing the ambiguity-set sample (and any stochastic optimization
randomness) fresh each trial while holding the true input and the observed
compressed update fixed. Then the probability of at least one success across
$K$ trials is
\begin{equation}
 P(K)=1-(1-p)^K,
 \label{eq:cumulative}
\end{equation}
and the number of trials needed to reach a target compromise probability
$\tau$ is $K(\tau)=\lceil \ln(1-\tau)/\ln(1-p)\rceil$. Using the measured
single-trial rate $p=0.04$ from Section~3.2, Eq.~\eqref{eq:cumulative} gives
$K(0.5){\approx}17$, $K(0.9){\approx}57$, and $K(0.99){\approx}113$
(Figure~\ref{fig:lucky}). None of these trial counts require additional
communication with the client or auxiliary data; they only require offline
compute, which is cheap relative to what an FL server already spends on
aggregation and training.

\begin{figure}[t]
  \centering
  \includegraphics[width=0.85\columnwidth]{figures/lucky_draw_curve}
  \caption{Cumulative compromise probability $P(K)=1-(1-p)^K$ for the
  measured single-trial success rate $p=0.04$ (Section~3.2). A defense that
  looks effective against one attack attempt is compromised with 90\%
  probability within about 56 independent, offline attack attempts.}
  \label{fig:lucky}
\end{figure}

We treat $p=0.04$ as a single operating point measured on one setting
(random sparsification imputation, Setup B) rather than a universal constant;
$p$ should depend on compression severity, batch size, and model, and our
independence assumption is optimistic because all trials condition on the
same observed compressed update and true input. We report
Eq.~\eqref{eq:cumulative} as a first-order model motivating the restart
strategy, not as a validated end-to-end guarantee; \todo{directly measure
$P(K)$ empirically by running $K$ independent trials per client/setting and
counting distinct successes, and fit $p$ as a function of compression level.}

\section{CoRA: consistency-ranked restart attack}
Section~3 identifies two requirements for turning repeated trials into a
reliable attack rather than a lucky occasional one: (i) candidates should be
generated from a noise/imputation model that respects the true, non-uniform,
per-tensor structure of the update instead of i.i.d.\ resampling, and
(ii) the adversary should \emph{select} a good candidate rather than average
across all of them, without folding the selection criterion into the
optimizer's loss (which over-constrains it, Section~3.3). CoRA implements
both, on top of any base analytical or optimization-based attack $A$ (DLG or
CPA/FastICA).

\textbf{Calibrated candidate generation.} Instead of i.i.d.\ noise over a
fixed range, we fit a per-tensor, per-layer noise/imputation distribution
$p_\theta(\cdot\mid\g_{\mathrm{obs}},\C)$ before the attack, using only
information already available to the adversary from the compression scheme
itself: the per-tensor quantization scale (for quantized updates) or the
empirical magnitude distribution of the retained coordinates in the same
tensor, truncated below the top-$k$ threshold (for sparsified updates).
Sampling missing/quantized values from this calibrated distribution, rather
than a shared i.i.d.\ range, directly targets the scale-mismatch failure
mode observed in Section~3.3.

\textbf{Independent unconstrained restarts.} For each of $K$ trials we draw
one calibrated candidate update $\tilde\g^{(t)}\sim p_\theta$ and run the
base attack $A$ unmodified: $\hat\x^{(t)}=A(\tilde\g^{(t)})$. We deliberately
do not add a compressed-gradient consistency term to $A$'s objective, since
Section~3.3 showed that doing so over-constrains the optimizer.

\textbf{Post-hoc consistency ranking.} For each candidate reconstruction, we
recompute its gradient $\g(\hat\x^{(t)})$ by a forward/backward pass through
the known model and loss, apply the \emph{same} compression operator the
client used, and score the candidate by how close this matches the actually
observed update:
\begin{equation}
 s_t=-\big\|\C\big(\g(\hat\x^{(t)})\big)-\g_{\mathrm{obs}}\big\|.
 \label{eq:score}
\end{equation}
This score is self-supervised: it uses only quantities the adversary already
has (the model, the compression scheme, and the observed compressed update),
never the true image. We return the top-scoring candidate,
$\hat\x^{(t^\ast)}$ with $t^\ast=\arg\max_t s_t$ (or a short list of the top-$m$
for a human analyst), instead of averaging over $t$.

\begin{algorithm}[t]
\caption{CoRA for one client's compressed update}
\label{alg:cora}
\begin{algorithmic}[1]
\REQUIRE observed compressed update $\g_{\mathrm{obs}}$, compression operator
$\C$, calibrated candidate model $p_\theta$, base attack $A$, trial budget $K$
\FOR{$t=1,\dots,K$}
  \STATE sample $\tilde\g^{(t)}\sim p_\theta(\cdot\mid\g_{\mathrm{obs}},\C)$
  \STATE $\hat\x^{(t)}\leftarrow A(\tilde\g^{(t)})$ \COMMENT{unconstrained base attack}
  \STATE compute $s_t$ from Eq.~\eqref{eq:score}
\ENDFOR
\STATE $t^\ast\leftarrow\arg\max_t s_t$
\RETURN $\hat\x^{(t^\ast)}$ (or the top-$m$ candidates by $s_t$)
\end{algorithmic}
\end{algorithm}

CoRA is attack-agnostic (it wraps DLG or CPA/FastICA without modifying their
internals), requires no auxiliary dataset or pretrained generative prior
(unlike attacks that strengthen the optimization prior itself
\cite{li2022ggl}), and is embarrassingly parallel across trials.
\todo{Implement and evaluate CoRA end-to-end against Setups A and B; report
success rate at fixed $K$ and reconstruction quality of the selected
candidate vs. (i) a single naive-decoding trial and (ii) naive $K$-trial
pixel averaging, at matched compute budget.}

\section{Experiments}
\subsection{Research questions}
(Q1) At what compression severity does single-trial reconstruction quality
collapse, and is the collapse sharp or gradual (Section~3.1)? (Q2) What is
the true per-trial success rate of resampling-based attacks, and does the
closed-form model of Eq.~\eqref{eq:cumulative} predict cumulative compromise
accurately (Section~3.4)? (Q3) Does CoRA's calibrated-sampling-plus-selection
design recover more successes per unit compute than naive i.i.d.
resampling-plus-averaging, and does it close the gap to the FP32 baseline
(Section~4)?

\subsection{Setup}
We reuse Setups A and B from Section~3.1 for all experiments. Baselines are
DLG \cite{zhu2019deep}, Inverting Gradients \cite{geiping2020inverting}, and
CPA/FastICA \cite{kariyappa2023cocktail}, each evaluated (i) once with naive
deterministic decoding, (ii) as an ensemble of $K$ i.i.d.-resampled,
pixel-averaged trials, and (iii) wrapped in CoRA. We report PSNR, SSIM,
LPIPS, and (for CPA) ICA/source cosine similarity, together with the
per-setting empirical success rate $p$ used to validate
Eq.~\eqref{eq:cumulative}. \todo{Run all three variants at
$K\in\{1,5,10,25,50\}$, 4-bit and 2-bit quantization, and global/layerwise
top-$k\in\{50,20,10\}\%$, over at least 100 examples and 3 seeds; add CIFAR-10
and \todo{one additional dataset relevant to the final ICASSP track}.}

Naive single-trial decoding and naive $K$-trial pixel averaging are already
reported in Tables~\ref{tab:fc2} and~\ref{tab:cpa}. Table~\ref{tab:planned}
adds the planned third arm, CoRA, at the same trial budget $K$.

\begin{table}[t]
\caption{Planned CoRA comparison at matched trial budget $K$, same settings
as Tables~\ref{tab:fc2}--\ref{tab:cpa}. All entries are placeholders pending
the experiments in Section~5.2.}
\label{tab:planned}
\centering
\small
\begin{tabular}{lccc}
\toprule
Method & Succ.\ rate & PSNR$\uparrow$ & SSIM$\uparrow$\\
\bottomrule
\end{tabular}
\end{table}

\subsection{Discussion}
\todo{Insert measured results once CoRA is implemented.} The central claim
to test is whether calibrated sampling and consistency-based selection
recover a substantially higher fraction of successes at fixed $K$ than naive
i.i.d.\ averaging, and whether the empirically measured $P(K)$ tracks
Eq.~\eqref{eq:cumulative} closely enough to be useful for auditing a
defense's real risk, rather than only its single-trial appearance.

\section{Related work}
Optimization-based gradient inversion reconstructs inputs by matching
observed and synthetic gradients \cite{zhu2019deep,geiping2020inverting}.
CPA instead treats fully-connected-layer gradients as a blind source
separation problem solved by ICA, and scales to large batches
\cite{kariyappa2023cocktail}. Gradient compression (quantization,
sparsification, and related obfuscation) is commonly proposed or assumed as
a side-effect privacy defense; Yue et al.~\cite{yue2023gradient} showed that
gradient \emph{obfuscation} more broadly can give a false sense of security
relative to differential privacy, though their analysis targets a different
compression regime and does not consider repeated stochastic attack trials.
Li et al.~\cite{li2022ggl} showed that defenses evaluated only against
relatively weak attacks can be broken by a \emph{stronger} attack that adds a
pretrained generative prior; our approach is complementary and orthogonal:
we do not strengthen the attack's prior, we exploit the fact that
information loss under compression is stochastic from the adversary's
viewpoint and combine repeated, calibrated resampling with self-supervised
selection. Zhang et al.~\cite{zhang2022survey} survey gradient inversion attacks and
defenses broadly, including compression-based defenses, but do not model
multi-trial compromise probability. Recent work has also proposed improved
attacks specifically targeting compressed gradients (e.g., quantization-aware
reconstruction losses) \cite{yu2024compressed}; unlike these, CoRA does not
require modifying the base attack's optimization objective (Section~3.3
shows that adding a compression-consistency term as a hard constraint
hurts), and instead separates candidate generation from selection.

\section{Conclusion}
Quantization and sparsification measurably degrade single-trial gradient
inversion attacks, but the degradation is closer to adding noise to a
lottery than to closing off the attack: a fixed, small per-trial success
probability, repeated over dozens of cheap, independent, offline attempts,
already yields a high chance of at least one successful reconstruction. A
naive way of running many trials (i.i.d.\ resampling and averaging) mostly
fails to realize this risk in practice; our proposed CoRA attack targets the
two reasons why, via calibrated candidate generation and post-hoc,
self-supervised selection instead of averaging.
\todo{State the final measured boundary: at what compression levels and
trial budgets does CoRA make compromise practical, and which of quantization
vs.\ sparsification is more resistant at matched communication savings.} We
recommend that defenses based on communication compression be evaluated by
their compromise probability under a realistic, repeated-trial adversary,
not by the outcome of a single deterministic attack attempt.

\bibliographystyle{IEEEbib}
\bibliography{refs}
\end{document}
